\documentclass[11pt]{article}

\usepackage[T1]{fontenc}
\usepackage[utf8]{inputenc}
\usepackage{lmodern}
\usepackage{microtype}
\usepackage[margin=1in]{geometry}
\usepackage{amsmath}
\usepackage{amssymb}
\usepackage{amsthm}
\usepackage{booktabs}
\usepackage{array}
\usepackage{enumitem}
\usepackage{xcolor}
\usepackage{fancyvrb}
\usepackage{url}
\usepackage[
  colorlinks=true,
  linkcolor=blue!55!black,
  citecolor=blue!55!black,
  urlcolor=blue!55!black,
  pdftitle={Six Certified Branch Exclusions for the Ternary Covering Problem K3(6,2)},
  pdfauthor={Ruturaj R Raval},
  pdfsubject={Exact computational coding theory and finite covering codes},
  pdfkeywords={covering codes, ternary codes, set cover duality, coding theory, computational combinatorics}
]{hyperref}

\newtheorem{theorem}{Theorem}
\newtheorem{lemma}[theorem]{Lemma}
\newtheorem{corollary}[theorem]{Corollary}
\newtheorem{remark}[theorem]{Remark}

\newcommand{\Hspace}{\mathcal{H}}
\newcommand{\dist}{d_{\mathrm H}}
\newcommand{\wt}{\operatorname{wt}}
\newcommand{\codeword}[1]{\texttt{#1}}

\title{Six Certified Branch Exclusions\\
for the Ternary Covering Problem \(K_3(6,2)\)}
\author{%
  Ruturaj R Raval\\
  \small Independent Researcher\\
  \small ORCID:
  \href{https://orcid.org/0000-0003-4930-8981}
       {\texttt{0000-0003-4930-8981}}
}
\date{September 6, 2026}

\begin{document}

\maketitle

\begin{abstract}
Let \(K_3(6,2)\) be the minimum number of radius-two Hamming balls needed
to cover the ternary space \(\{0,1,2\}^6\).  The published interval remains
\(15\leq K_3(6,2)\leq 17\).  A hypothetical cover with at most sixteen
centers can be translated to contain \(\codeword{000000}\).  An antipodal
sphere argument excludes normalized maximum weight at most four, and a
radius-three sphere argument forces an additional center of weight at most
four.  Stabilizer orbits then give 24 normalized third-center branches for
maximum weight five and 14 for maximum weight six.

This report certifies six of those 38 branches as impossible.  For each
branch, nonnegative integer weights are assigned to the points not covered
by its three fixed centers.  Every remaining admissible center covers
weighted capacity at most \(Q\), while the total weight \(W\) satisfies
\(W>13Q\).  After scaling by \(Q\), these are exact feasible dual solutions
for the residual fractional set-cover relaxations with objective greater
than thirteen.  Five weight-five branches and one weight-six branch are
therefore excluded.  Independent Python and C++20 verifiers reconstruct
the complete branch definitions, evaluate all 729 ambient words and every
admissible center, and agree on every exact count.  The normalized frontier
is reduced from 38 branches to 32.  No new global bound is claimed.
\end{abstract}

\section{Introduction}

A \(q\)-ary covering code of length \(n\) and radius \(R\) is a subset of
\(\{0,\ldots,q-1\}^n\) whose radius-\(R\) Hamming balls cover the entire
space.  The minimum possible cardinality is denoted by \(K_q(n,R)\).
Covering codes connect coding theory, finite geometry, domination,
set cover, and exact optimization.

This report studies the unresolved value \(K_3(6,2)\).  The historical
covering-code table records
\[
  15\leq K_3(6,2)\leq 17
\]
\cite{keri-table}.  The upper bound is attributed to work of
H\"am\"al\"ainen and Rankinen from 1991
\cite{hamalainen-rankinen-1991}; the lower bound is recorded in the updated
binary and ternary mixed-code table of Bertolo, \"Osterg{\aa}rd, and Weakley
from 2004 \cite{bertolo-ostergard-weakley-2004}.  A modern Lean
formalization supplies an explicit checked 17-word construction
\cite{covering-codes-lean-upper}.

Recent semidefinite methods improve many covering-code lower bounds, but
their numerical value for this cell remains below the established integer
lower bound fifteen \cite{gijswijt-polak-2026}.  Public exploratory notes
for the exact 16-center case report feasible anchored relaxations and
third-orbit residual LP scans that pruned none of their tested branches
\cite{covering-codes-lean-failures}.  They identify verified finite search
and proof-producing branch decompositions as natural next directions.

The contribution here is a scoped exact result.  The normalized
size-at-most-16 problem is first reduced to 38 third-center branches.  Six
branches then receive short residual set-cover dual certificates.  The
distinction between this result and a complete solution is maintained
throughout:

\begin{itemize}[leftmargin=2em]
  \item six named normalized branches are rigorously excluded;
  \item 32 normalized third-center branches remain unresolved;
  \item no 16-word covering code is constructed;
  \item no complete exclusion of size-at-most-16 covers is proved; and
  \item the global interval \(15\leq K_3(6,2)\leq 17\) is unchanged.
\end{itemize}

\section{Definitions}

Let
\[
  \Hspace=\{0,1,2\}^6.
\]
For \(x,y\in\Hspace\), write
\[
  \dist(x,y)=|\{i:x_i\neq y_i\}|
\]
for Hamming distance and
\[
  B_2(c)=\{x\in\Hspace:\dist(x,c)\leq 2\}
\]
for the radius-two ball centered at \(c\).  A set \(C\subseteq\Hspace\)
is a radius-two covering code when
\[
  \bigcup_{c\in C}B_2(c)=\Hspace.
\]

The ambient space contains \(3^6=729\) words, and every radius-two ball
contains
\[
  1+\binom61 2+\binom62 2^2=73
\]
words.

For \(x=(x_1,\ldots,x_6)\), define its weight by
\[
  \wt(x)=|\{i:x_i\neq 0\}|.
\]
Translation by a word of \(\Hspace\), coordinate permutations, and
independent symbol permutations in each coordinate preserve Hamming
distance and covering.

\section{The normalized 38-branch frontier}

\subsection{Excluding maximum weight at most four}

\begin{lemma}[Antipodal sphere]
\label{lem:antipodal}
In a radius-two cover with at most sixteen centers, every selected center
has another selected center at distance at least five.
\end{lemma}

\begin{proof}
Fix a selected center \(z\).  The sphere of radius six around \(z\) has
64 words.  A radius-two ball whose center is at distance four, five, or six
from \(z\) covers respectively 4, 12, or 22 words of that sphere.  A center
at smaller distance covers none.

If all other centers were at distance at most four from \(z\), at most
fifteen balls would contribute, each with capacity at most four.  Their
total capacity would be at most \(15\cdot4=60<64\), a contradiction.
\end{proof}

Translate a selected center to \(\codeword{000000}\).  Lemma
\ref{lem:antipodal} implies that the translated code contains a word of
weight at least five.  Choosing a maximum-weight word and normalizing its
nonzero symbols gives the anchor
\[
  a_d=1^d0^{6-d}
\]
with \(d\in\{5,6\}\).  Thus the normalized maximum-weight-at-most-four
branch is impossible.

\subsection{A low-weight third center}

\begin{lemma}[Radius-three sphere]
\label{lem:radius-three}
In a radius-two cover with at most sixteen centers, every selected center
has another selected center at distance at most four.
\end{lemma}

\begin{proof}
Fix a selected center \(z\).  Its radius-three sphere has 160 words.  A
radius-two ball centered at distances five and six from \(z\) covers at
most 10 and 0 of those words, respectively.  If all other centers were at
distance at least five, their total capacity would be at most
\[
  15\cdot10=150<160,
\]
a contradiction.
\end{proof}

Apply Lemma \ref{lem:radius-three} to \(\codeword{000000}\).  Besides the
anchor of weight five or six, the code contains a third center of weight at
most four.

Fix \(d\in\{5,6\}\) and the pair
\(\{\codeword{000000},a_d\}\).  Its pointwise stabilizer permutes the first
\(d\) coordinates, permutes the remaining coordinates, and swaps symbols
1 and 2 independently outside the anchor support.  The orbit of a candidate
third center is determined by:

\begin{itemize}[leftmargin=2em]
  \item its counts of symbols 0, 1, and 2 in the first \(d\) coordinates;
  \item its number of nonzero symbols in the final \(6-d\) coordinates.
\end{itemize}

Restricting only the canonical third center to weight at most four gives
24 orbits for \(d=5\) and 14 orbits for \(d=6\).  Choosing the earliest
occupied orbit and mapping one selected member to its canonical
representative gives a complete nonoverlapping 38-branch split.  Fourth and
later centers retain every weight permitted by the anchor.

\section{Residual set-cover dual certificates}

Fix one normalized branch.  Let \(t\) be its third-center representative
and let \(E\) be the union of all earlier third-center orbits.  Define
\[
  F=\{\codeword{000000},a_d,t\}
\]
and
\[
  U=\{c\in\Hspace:\wt(c)\leq d\}\setminus(F\cup E).
\]
The set \(U\) includes every remaining center permitted by the branch:
unused members of the selected orbit, all later low-weight orbits, and every
heavier center allowed by the maximum-weight anchor.

Let
\[
  H=\{p\in\Hspace:\dist(p,f)>2\text{ for every }f\in F\}
\]
be the holes left by the three fixed centers.

\begin{theorem}[Weighted residual dual certificate]
\label{thm:dual}
Suppose there are nonnegative weights
\(w:H\rightarrow\mathbb{Z}_{\geq0}\) and an integer \(Q>0\) such that
\[
  K(c):=
  \sum_{\substack{p\in H\\\dist(p,c)\leq2}}w(p)
  \leq Q
  \qquad\text{for every }c\in U,
\]
while
\[
  W:=\sum_{p\in H}w(p)>13Q.
\]
Then the branch contains no radius-two covering code with at most sixteen
centers.
\end{theorem}

\begin{proof}
Any branch cover contains the three fixed centers and at most thirteen
additional centers \(T\subseteq U\).  Every \(p\in H\) must be covered by
at least one member of \(T\).  Since all weights are nonnegative,
\[
  W
  \leq
  \sum_{c\in T}
  \sum_{\substack{p\in H\\\dist(p,c)\leq2}}w(p)
  =
  \sum_{c\in T}K(c)
  \leq
  |T|Q
  \leq
  13Q,
\]
contradicting \(W>13Q\).
\end{proof}

\begin{remark}
After scaling by \(Q\), the values \(w(p)/Q\) form a feasible dual solution
for the residual fractional set-cover relaxation.  The objective value is
\(W/Q>13\), so the residual branch needs more than thirteen centers even
fractionally.
\end{remark}

\section{Six exact certificates}

The six certified branches and their exact values are listed in Table
\ref{tab:certificates}.

\begin{table}[htbp]
\centering
\caption{Exact residual dual certificates.}
\label{tab:certificates}
\begin{tabular}{rrrrrrrr}
\toprule
Anchor & Orbit & Representative & \(|U|\) & \(|H|\) & \(W\) & \(Q\) & \(13Q\)\\
\midrule
5 & 19 & \codeword{011110} & 270 & 549 & 80 & 6 & 78\\
5 & 20 & \codeword{011120} & 265 & 541 & 40 & 3 & 39\\
5 & 21 & \codeword{011220} & 245 & 528 & 40 & 3 & 39\\
5 & 22 & \codeword{012220} & 215 & 522 & 40 & 3 & 39\\
5 & 23 & \codeword{022220} & 195 & 516 & 80 & 6 & 78\\
6 & 13 & \codeword{002222} & 269 & 516 & 132 & 10 & 130\\
\bottomrule
\end{tabular}
\end{table}

The five weight-five certificates all have dual objective
\[
  \frac{W}{Q}=\frac{40}{3}>13.
\]
The weight-six certificate has objective
\[
  \frac{W}{Q}=\frac{66}{5}>13.
\]
Theorem \ref{thm:dual} therefore excludes all six branches.

\subsection{Compact orbit notation}

For a weight-five branch, partition the first five coordinates according to
whether \(t_i\) is 0, 1, or 2.  For each group, record the counts of symbols
0, 1, and 2 in the ambient point.  These three count blocks are followed by
a bit indicating whether coordinate 6 is zero or nonzero.  Thus a key such
as
\[
  \codeword{001|202|000|0}
\]
describes one orbit of ambient points under the stabilizer of the three
fixed centers.

For the weight-six branch \(t=\codeword{002222}\), the two blocks record
symbol counts in the two coordinates where \(t_i=0\) and the four
coordinates where \(t_i=2\).

Appendix \ref{app:weights} lists every positive-weight orbit.  Unlisted holes
have weight zero.

\section{Independent exact verification}

Two implementations verify the certificates:

\begin{itemize}[leftmargin=2em]
  \item a Python checker reads the machine-readable JSON certificate file;
  \item an independent C++20 checker contains a separately encoded
        certificate table and reconstructs the geometry independently.
\end{itemize}

For every branch, each checker:

\begin{enumerate}[leftmargin=2em]
  \item reconstructs the complete weight-at-most-four third-center orbit
        partition;
  \item validates the branch representative and all earlier forbidden
        orbits;
  \item retains every remaining center of weight at most the anchor weight;
  \item reconstructs all holes left by the three fixed centers;
  \item expands every positive-weight point orbit and checks its size;
  \item evaluates \(W\) and \(K(c)\) for every individual \(c\in U\); and
  \item checks the strict integer inequality \(W>13Q\).
\end{enumerate}

The implementations produce byte-identical summaries:

\begin{Verbatim}[fontsize=\small]
w5 orbit 19 011110: |U|=270 |H|=549 W=80 Q=6 13Q=78 margin=2
w5 orbit 20 011120: |U|=265 |H|=541 W=40 Q=3 13Q=39 margin=1
w5 orbit 21 011220: |U|=245 |H|=528 W=40 Q=3 13Q=39 margin=1
w5 orbit 22 012220: |U|=215 |H|=522 W=40 Q=3 13Q=39 margin=1
w5 orbit 23 022220: |U|=195 |H|=516 W=80 Q=6 13Q=78 margin=2
w6 orbit 13 002222: |U|=269 |H|=516 W=132 Q=10 13Q=130 margin=2
all weighted branch certificates verified
\end{Verbatim}

A second canonical dump requires the JSON certificate data and the separately
encoded C++20 table to agree entry for entry.  For all six branches, the
fixed, forbidden, and admissible center sets reconstructed by the Python
checker are also compared with manifests emitted by the production CNF
generator.

Mutation tests alter representatives, point weights, numeric types, orbit
sizes, expected totals, capacity bounds, and key uniqueness.  Every malformed
certificate produced by these mutation classes is rejected.

The primary replay commands are:

\begin{Verbatim}[fontsize=\small]
make
python3 tools/verify_branch_certificates.py
build/verify_branch_certificates
make test
\end{Verbatim}

\section{Significance and limitations}

The result supplies exact local lower bounds inside a complete normalized
partition.  Within each certified branch, the three fixed centers require at
least fourteen additional centers, so any cover in that branch has at least
seventeen centers.

The certificates are reusable in deeper search.  A branch-and-bound, SAT,
MILP, or exact-cover implementation may reject any node whose normalized
ancestor is one of the six certified branches.  The small integer dual
weights also provide regression instances for set-cover preprocessing and
symmetry handling.

The limitations are equally explicit:

\begin{itemize}[leftmargin=2em]
  \item the six branches are only part of the normalized 38-branch frontier;
  \item 19 weight-five branches and 13 weight-six branches remain unresolved;
  \item the certificates do not rule out a 16-word cover in another branch;
  \item no 15-word or 16-word cover is constructed; and
  \item the exact value of \(K_3(6,2)\) remains open.
\end{itemize}

\section{Conclusion}

Elementary sphere-capacity arguments reduce the normalized size-at-most-16
problem for \(K_3(6,2)\) to 38 third-center branches.  Exact residual
set-cover dual certificates exclude six of those branches, with independent
Python and C++20 replay over the full ternary space.  The certified
normalized frontier now contains 32 unresolved branches.

The next finite step is to apply deeper stabilizer-orbit splitting to those
32 branches and search for additional exact dual certificates.  Independent
construction search remains essential because one verified 16-word cover
would settle the upper-bound direction immediately.

\appendix

\section{Positive-weight point orbits}
\label{app:weights}

Each entry below has the form
\[
  \texttt{orbit-key : weight (orbit-size)}.
\]

\subsection*{\(\codeword{011110}\)}
\begin{Verbatim}[fontsize=\small]
001|202|000|0 : 1 (6)
001|310|000|1 : 2 (8)
010|202|000|0 : 1 (6)
010|310|000|1 : 2 (8)
100|202|000|1 : 1 (12)
100|211|000|1 : 1 (24)
\end{Verbatim}

\subsection*{\(\codeword{011120}\)}
\begin{Verbatim}[fontsize=\small]
001|300|001|1 : 2 (2)
001|300|010|1 : 2 (2)
010|300|001|1 : 2 (2)
010|300|010|1 : 2 (2)
100|102|100|1 : 1 (6)
100|111|100|1 : 1 (12)
100|120|100|1 : 1 (6)
\end{Verbatim}

\subsection*{\(\codeword{011220}\)}
\begin{Verbatim}[fontsize=\small]
001|101|200|1 : 1 (4)
001|110|200|1 : 1 (4)
010|101|200|1 : 1 (4)
010|110|200|1 : 1 (4)
100|002|200|1 : 1 (2)
100|011|200|1 : 1 (4)
100|020|200|1 : 1 (2)
100|200|002|1 : 2 (2)
100|200|011|1 : 2 (4)
100|200|020|1 : 2 (2)
\end{Verbatim}

\subsection*{\(\codeword{012220}\)}
\begin{Verbatim}[fontsize=\small]
001|001|300|1 : 2 (2)
001|010|300|1 : 2 (2)
010|001|300|1 : 2 (2)
010|010|300|1 : 2 (2)
100|100|102|1 : 1 (6)
100|100|111|1 : 1 (12)
100|100|120|1 : 1 (6)
\end{Verbatim}

\subsection*{\(\codeword{022220}\)}
\begin{Verbatim}[fontsize=\small]
001|000|211|0 : 1 (12)
001|000|301|1 : 2 (8)
010|000|211|0 : 1 (12)
010|000|301|1 : 2 (8)
100|000|121|0 : 1 (12)
100|000|220|1 : 1 (12)
\end{Verbatim}

\subsection*{\(\codeword{002222}\)}
\begin{Verbatim}[fontsize=\small]
002|301 : 1 (4)
002|310 : 1 (4)
011|310 : 1 (8)
020|301 : 1 (4)
101|202 : 1 (12)
101|211 : 1 (24)
101|220 : 1 (12)
110|211 : 2 (24)
200|121 : 1 (12)
200|130 : 1 (4)
\end{Verbatim}

\begin{thebibliography}{9}

\bibitem{keri-table}
G.~K\'eri,
\emph{Tables for covering codes},
online table,
\url{https://old.sztaki.hu/~keri/codes/3_tables.pdf},
accessed September 6, 2026.

\bibitem{hamalainen-rankinen-1991}
H.~H\"am\"al\"ainen and S.~Rankinen,
\emph{Upper bounds for football pool problems and mixed covering codes},
Journal of Combinatorial Theory, Series A 56 (1991), 84--95.
\url{https://doi.org/10.1016/0097-3165(91)90024-B}.

\bibitem{bertolo-ostergard-weakley-2004}
R.~Bertolo, P.~R.~J.~\"Osterg{\aa}rd, and W.~D.~Weakley,
\emph{An updated table of binary/ternary mixed covering codes},
Journal of Combinatorial Designs 12 (2004), 157--176.
\url{https://doi.org/10.1002/jcd.20008}.

\bibitem{gijswijt-polak-2026}
D.~Gijswijt and S.~Polak,
\emph{Semidefinite lower bounds for covering codes},
arXiv:2504.01932v2, 2026.
\url{https://arxiv.org/abs/2504.01932}.

\bibitem{covering-codes-lean-upper}
A.~Florath,
\emph{covering-codes-lean: explicit upper bound for \(K_3(6,2)\)},
file \path{CoveringCodes/Database/Sources/SmallExplicitUpper/K_3_6_2.lean},
GitHub repository, 2026.
\url{https://github.com/florath/covering-codes-lean}.

\bibitem{covering-codes-lean-failures}
A.~Florath,
\emph{Failed attempts for \(K_3(6,2)\)},
file \path{docs/failures/K_3_6_2.md},
covering-codes-lean project notes, 2026.
\url{https://github.com/florath/covering-codes-lean}.

\end{thebibliography}

\end{document}
